\chapter{Precalculus Reference}

\textit{This appendix collects the precalculus facts and identities
used throughout the course. Results are stated without proof.
Familiarity with this material is assumed from the first chapter onward.}

%% ─── SECTION A.1 ─────────────────────────────────────────────────────────────
\section{Trigonometric Identities}
\label{sec:A.1}

Trigonometric identities pervade every part of calculus. When you
differentiate $\sin x$ or $\cos x$, when you integrate powers of
trigonometric functions, and when you simplify expressions after a
trigonometric substitution, you rely on the identities collected here.
Knowing them cold---not just recognising them, but being able to
reproduce them quickly---will save you enormous time throughout the
course.

Many of the identities below are consequences of just one fact: the
Pythagorean identity $\sin^2\theta + \cos^2\theta = 1$.  We organise
them into groups so that you can see the logical relationships.

\begin{proposition}[Pythagorean Identities]\label{prop:A-1}
\begin{align*}
  \sin^2\theta + \cos^2\theta &= 1 \\
  1 + \tan^2\theta            &= \sec^2\theta \\
  1 + \cot^2\theta            &= \csc^2\theta
\end{align*}
\end{proposition}

\begin{proposition}[Reciprocal and Quotient Identities]\label{prop:A-2}
\begin{align*}
  \csc\theta &= \dfrac{1}{\sin\theta} \\
  \sec\theta &= \dfrac{1}{\cos\theta} \\
  \cot\theta &= \dfrac{1}{\tan\theta} \\
  \tan\theta &= \dfrac{\sin\theta}{\cos\theta} \\
  \cot\theta &= \dfrac{\cos\theta}{\sin\theta}
\end{align*}
\end{proposition}

\begin{proposition}[Angle Addition and Subtraction Formulas]\label{prop:A-3}
\begin{align*}
  \sin(\alpha \pm \beta) &= \sin\alpha\cos\beta \pm \cos\alpha\sin\beta \\
  \cos(\alpha \pm \beta) &= \cos\alpha\cos\beta \mp \sin\alpha\sin\beta \\
  \tan(\alpha \pm \beta) &= \dfrac{\tan\alpha \pm \tan\beta}
                                  {1 \mp \tan\alpha\tan\beta}
\end{align*}
\end{proposition}

\begin{proposition}[Double Angle Formulas]\label{prop:A-4}
\begin{align*}
  \sin(2\theta) &= 2\sin\theta\cos\theta \\
  \cos(2\theta) &= \cos^2\theta - \sin^2\theta \\
                &= 2\cos^2\theta - 1 \\
                &= 1 - 2\sin^2\theta \\
  \tan(2\theta) &= \dfrac{2\tan\theta}{1 - \tan^2\theta}
\end{align*}
\end{proposition}

\begin{proposition}[Half Angle Formulas]\label{prop:A-5}
\begin{align*}
  \sin^2\theta &= \dfrac{1 - \cos(2\theta)}{2} \\
  \cos^2\theta &= \dfrac{1 + \cos(2\theta)}{2} \\
  \tan^2\theta &= \dfrac{1 - \cos(2\theta)}{1 + \cos(2\theta)}
\end{align*}
\end{proposition}

\begin{proposition}[Product-to-Sum Formulas]\label{prop:A-6}
\begin{align*}
  \sin\alpha\cos\beta &= \tfrac{1}{2}[\sin(\alpha+\beta)+\sin(\alpha-\beta)] \\
  \cos\alpha\sin\beta &= \tfrac{1}{2}[\sin(\alpha+\beta)-\sin(\alpha-\beta)] \\
  \cos\alpha\cos\beta &= \tfrac{1}{2}[\cos(\alpha-\beta)+\cos(\alpha+\beta)] \\
  \sin\alpha\sin\beta &= \tfrac{1}{2}[\cos(\alpha-\beta)-\cos(\alpha+\beta)]
\end{align*}
\end{proposition}

\begin{proposition}[Sum-to-Product Formulas]\label{prop:A-7}
\begin{align*}
  \sin\alpha + \sin\beta
    &= 2\sin\!\left(\dfrac{\alpha+\beta}{2}\right)
        \cos\!\left(\dfrac{\alpha-\beta}{2}\right) \\
  \sin\alpha - \sin\beta
    &= 2\cos\!\left(\dfrac{\alpha+\beta}{2}\right)
        \sin\!\left(\dfrac{\alpha-\beta}{2}\right) \\
  \cos\alpha + \cos\beta
    &= 2\cos\!\left(\dfrac{\alpha+\beta}{2}\right)
        \cos\!\left(\dfrac{\alpha-\beta}{2}\right) \\
  \cos\alpha - \cos\beta
    &= -2\sin\!\left(\dfrac{\alpha+\beta}{2}\right)
         \sin\!\left(\dfrac{\alpha-\beta}{2}\right)
\end{align*}
\end{proposition}

\begin{table}[h]
\centering
\caption{Values of Trigonometric Functions at Standard Angles}
\begin{tabular}{@{}ccccccc@{}}
\toprule
$\theta$
  & $\sin\theta$ & $\cos\theta$ & $\tan\theta$
  & $\csc\theta$ & $\sec\theta$ & $\cot\theta$ \\
\midrule
$0$
  & $0$ & $1$ & $0$
  & $\text{undef}$ & $1$ & $\text{undef}$ \\
$\pi/6$
  & $1/2$ & $\sqrt{3}/2$ & $1/\sqrt{3}$
  & $2$ & $2/\sqrt{3}$ & $\sqrt{3}$ \\
$\pi/4$
  & $\sqrt{2}/2$ & $\sqrt{2}/2$ & $1$
  & $\sqrt{2}$ & $\sqrt{2}$ & $1$ \\
$\pi/3$
  & $\sqrt{3}/2$ & $1/2$ & $\sqrt{3}$
  & $2/\sqrt{3}$ & $2$ & $1/\sqrt{3}$ \\
$\pi/2$
  & $1$ & $0$ & $\text{undef}$
  & $1$ & $\text{undef}$ & $0$ \\
$\pi$
  & $0$ & $-1$ & $0$
  & $\text{undef}$ & $-1$ & $\text{undef}$ \\
$3\pi/2$
  & $-1$ & $0$ & $\text{undef}$
  & $-1$ & $\text{undef}$ & $0$ \\
$2\pi$
  & $0$ & $1$ & $0$
  & $\text{undef}$ & $1$ & $\text{undef}$ \\
\bottomrule
\end{tabular}
\end{table}

\begin{remark}
All angles in the table above are measured in radians. The notation
$\text{undef}$ indicates that the function is undefined at that angle
because its denominator is zero; geometrically, the function has a
vertical asymptote there.
\end{remark}

%% ─── SECTION A.2 ─────────────────────────────────────────────────────────────
\section{Exponential and Logarithmic Laws}
\label{sec:A.2}

The exponential function $e^x$ and the natural logarithm $\ln x$ are
everywhere in calculus: they appear in differentiation rules, integration
formulas, growth and decay models, and the definitions of many other
functions. The algebraic laws listed here are the foundation on which all
of that rests.

In particular, the change-of-base formula lets us reduce any exponential
$a^x$ to $e^{x \ln a}$, which is essential when differentiating or
integrating exponentials with arbitrary bases in Chapters~3 and~4.

\begin{proposition}[Laws of Exponents]\label{prop:A-8}
For $a > 0$ and all $x, y \in \R$:
\begin{align*}
  a^x \cdot a^y    &= a^{x+y} \\
  \dfrac{a^x}{a^y} &= a^{x-y} \\
  (a^x)^y          &= a^{xy} \\
  (ab)^x           &= a^x b^x \\
  a^0              &= 1 \\
  a^{-x}           &= \dfrac{1}{a^x}
\end{align*}
\end{proposition}

\begin{remark}
The law $(a^x)^y = a^{xy}$ is used constantly when simplifying
expressions before differentiating---for instance, rewriting
$\sqrt[3]{x^2}$ as $x^{2/3}$ so that the power rule applies directly.
\end{remark}

\begin{proposition}[Laws of Logarithms]\label{prop:A-9}
For $a > 0$, $a \neq 1$, and $x, y > 0$:
\begin{align*}
  \log_a(xy)  &= \log_a x + \log_a y \\
  \log_a(x/y) &= \log_a x - \log_a y \\
  \log_a(x^r) &= r\log_a x \\
  \log_a x    &= \dfrac{\ln x}{\ln a}
                 \quad \text{(change of base)} \\
  \ln(e^x)    &= x \quad \text{for all } x \in \R \\
  e^{\ln x}   &= x \quad \text{for all } x > 0
\end{align*}
\end{proposition}

\begin{remark}
The change-of-base formula $\log_a x = \dfrac{\ln x}{\ln a}$ is the key
to differentiating $a^x$ for $a \neq e$: we write $a^x = e^{x \ln a}$
and apply the chain rule. It also converts any logarithmic integral into
one involving $\ln$.
\end{remark}

%% ─── SECTION A.3 ─────────────────────────────────────────────────────────────
\section{Algebraic Techniques}
\label{sec:A.3}

The algebraic manipulations below appear throughout the course: completing
the square is essential for certain integrals and limit calculations,
factoring identities simplify rational expressions, and the binomial
theorem underpins Taylor expansions and series analysis.

\begin{proposition}[Completing the Square]\label{prop:A-10}
\[
  x^2 + bx = \left(x + \dfrac{b}{2}\right)^2 - \dfrac{b^2}{4}
\]
\end{proposition}

\begin{example}\label{ex:A-1}
Write $x^2 + 6x + 5$ in completed square form.

We isolate the quadratic and linear terms and apply the formula:
\begin{align*}
  x^2 + 6x + 5
    &= (x^2 + 6x) + 5 \\
    &= \left(x^2 + 6x + 9\right) - 9 + 5 \\
    &= (x + 3)^2 - 4
\end{align*}
Hence $x^2 + 6x + 5 = (x+3)^2 - 4$, so the minimum value of the
expression is $-4$, attained at $x = -3$.
\end{example}

\begin{proposition}[Factoring Identities]\label{prop:A-11}
\begin{align*}
  a^2 - b^2 &= (a-b)(a+b) \\
  a^3 - b^3 &= (a-b)(a^2+ab+b^2) \\
  a^3 + b^3 &= (a+b)(a^2-ab+b^2)
\end{align*}
\end{proposition}

\begin{proposition}[Binomial Theorem]\label{prop:A-12}
For $n \in \N$ and $a, b \in \R$:
\[
  (a+b)^n = \sum_{k=0}^{n} \binom{n}{k} a^{n-k} b^k
\]
The first three cases explicitly:
\begin{align*}
  (a+b)^2 &= a^2 + 2ab + b^2 \\
  (a+b)^3 &= a^3 + 3a^2 b + 3ab^2 + b^3
\end{align*}
\end{proposition}

\begin{remark}[Polynomial Long Division]
When integrating a rational function $p(x)/q(x)$ where
$\deg p \geq \deg q$, perform polynomial long division first to write
the integrand as a polynomial plus a proper rational function.

For example, divide $x^3 + 2x$ by $x^2 + 1$:
\begin{align*}
  x^3 + 2x &= (x^2 + 1) \cdot x + (2x - x) \\
            &= (x^2 + 1) \cdot x + x
\end{align*}

We verify this step by step. We seek $q(x)$ and $r(x)$ with
$\deg r < \deg(x^2 + 1) = 2$ such that $x^3 + 2x = (x^2+1)\,q(x) + r(x)$.

First, divide the leading term: $x^3 \div x^2 = x$.

Multiply and subtract:
\begin{align*}
  x^3 + 2x - x \cdot (x^2 + 1)
    &= x^3 + 2x - x^3 - x \\
    &= x
\end{align*}

Since $\deg(x) = 1 < 2 = \deg(x^2+1)$, we stop. Therefore:
\[
  \dfrac{x^3 + 2x}{x^2 + 1} = x + \dfrac{x}{x^2 + 1}
\]
\end{remark}

%% ─── SECTION A.4 ─────────────────────────────────────────────────────────────
\section{Inverse Functions}
\label{sec:A.4}

Inverse functions are central to Chapter~4, where we study inverse
trigonometric functions ($\arcsin$, $\arccos$, $\arctan$) and where we
treat the natural logarithm as the inverse of the exponential. The
derivative rule for inverse functions, stated below, appears in Chapter~3
and is used repeatedly whenever we differentiate an inverse function.

Understanding when a function \textit{has} an inverse---and the role of
restricting the domain to guarantee invertibility---is essential before
any of those later results make sense.

\begin{definition}[Inverse Function]\label{def:A-1}
Let $f \colon A \to B$ be a function. A function $g \colon B \to A$
is called the \emph{inverse} of $f$ if
\[
  g(f(x)) = x \quad \text{for all } x \in A
  \qquad \text{and} \qquad
  f(g(y)) = y \quad \text{for all } y \in B.
\]
When such a $g$ exists, $f$ is called \emph{invertible}, and we write
$g = f^{-1}$.
\end{definition}

\begin{remark}
A function $f$ is invertible if and only if it is \textit{bijective}:
injective (one-to-one) and surjective (onto).
Geometrically, $f$ passes the horizontal line test.
\end{remark}

\begin{proposition}[Inverse Function Derivative Rule]\label{prop:A-13}
Let $f$ be differentiable and invertible on an interval,
with $f'(x) \neq 0$. Then $f^{-1}$ is differentiable and
\[
  \bigl(f^{-1}\bigr)'(y) = \dfrac{1}{f'\!\left(f^{-1}(y)\right)}.
\]
\end{proposition}

\begin{remark}
The proof of \cref{prop:A-13} uses the chain rule and appears in
\cref{sec:4.4}.
\end{remark}
